\section{State of the Art and Baselines}
\label{sec:baselines_sota}

\paragraph{State of the Art in Representation Post-Processing}
Pretrained textual representations commonly concentrate within narrow cones, creating high average cosine similarities between unrelated sequences \citep{ethayarajh-2019-contextual}. The state of the art in training-free post-processing spans three major methodological paradigms:
\begin{itemize}
\item Coordinate Paradigm: Identifies outlier rogue dimensions that dominate inner products and applies coordinate-wise z-score standardization \citep{timkey2021all}.
\item Subspace Paradigm: Treats origin shift as a stable mean bias vector and applies orthogonal projection deflation (R2) to cancel parallel sample-mean estimation noise \citep{ren2025correcting}.
\item Spectral Paradigm: Regularizes covariance eigenvalues to enforce hyperspherical isotropy via eigenvalue-regularized Zero-phase Component Analysis (Soft-ZCA) \citep{diera2024isotropy} or adapts tempering strength based on signal-to-noise ratio knee points (SpecTemp) \citep{li2026spectral}.
\end{itemize}

\paragraph{Competitive Model Baselines}
We evaluate four open-weight models in two tiers, dedicated contrastive encoders and un-fine-tuned base causal language models (Table~\ref{tab:model_architectures}).

\begin{table*}[t]
\centering
\small
\caption{Evaluated models. Base models are evaluated under both mean pooling and final non-padding token pooling.}
\label{tab:model_architectures}
\resizebox{\textwidth}{!}{%
\begin{tabular}{>{\raggedright\arraybackslash}p{3.1cm}>{\raggedright\arraybackslash}p{1.7cm}cccc >{\raggedright\arraybackslash}p{1.8cm}>{\raggedright\arraybackslash}p{2.4cm}>{\raggedright\arraybackslash}p{2.6cm}}
\toprule
\textbf{Model} & \textbf{Tier} & \textbf{Params} & \textbf{Native $d$} & \textbf{Context} & \textbf{MRL} & \textbf{Pooling} & \textbf{Training Objective} & \textbf{Prompts and Output} \\
\midrule
\texttt{google/\allowbreak embeddinggemma-300m} \citep{vera-etal-2025-embeddinggemma} & Dedicated contrastive & 308M & 768 & 2,048 & Yes & Mean & Contrastive InfoNCE & Asymmetric instruction prompts; unit-normalized \\
\texttt{Qwen/\allowbreak Qwen3-Embedding-0.6B} \citep{zhang-etal-2025-qwen3embedding} & Dedicated contrastive & 600M & 1,024 & 32,768 & Yes & Final non-padding token & Contrastive InfoNCE & Asymmetric instruction prompts; unit-normalized \\
\texttt{google/\allowbreak gemma-3-1b-pt} \citep{gemmateam-2025-gemma3} & Base causal LM & 1.0B & 1,152 & 32,768 & No & Mean and final token & Next-token cross-entropy, tied embeddings & No prompts; no sentence-level ranking loss \\
\texttt{Qwen/\allowbreak Qwen3.5-0.8B-Base} \citep{qwenteam-2026-qwen35} & Base causal LM & 0.8B & 1,024 & 262,144 & No & Mean and final token & Next-token cross-entropy, tied embeddings & No prompts; no sentence-level ranking loss \\
\bottomrule
\end{tabular}%
}
\end{table*}

\paragraph{Availability of Baseline Implementations}
Model checkpoints and tokenizers are publicly available through the Hugging Face Hub. Matryoshka Representation Learning (MRL) follows \citet{kusupati2022matryoshka}. Regarding post-processing methods:
\begin{itemize}
\item Full whitening implementations follow public algorithms from \citet{su2021whitening} and \citet{huang2021whiteningbert}.
\item Soft-ZCA whitening follows the spectral formulation from \citet{diera2024isotropy}.
\item SpecTemp compression is implemented from the SNR knee detection equations and Kneedle algorithm specified by \citet{li2026spectral}.
\item R1 mean subtraction, R2 subspace deflation, ABTT component purging, and scientific controls (random direction deflation, unrenormalized centering) are implemented from the mathematical derivations in \citet{ren2025correcting} and \citet{mu2018all}, as no public repository was provided with the original publication.
\end{itemize}
All methods are implemented as an FP32 PyTorch codebase.
